% akilan.tex
\documentclass[11pt]{book}
\usepackage[utf8]{inputenc}
\usepackage[T1]{fontenc}
\usepackage{geometry}
\geometry{a5paper, margin=1in}
\usepackage{setspace}
\onehalfspacing
\usepackage{microtype}
\usepackage{ebgaramond}
\usepackage{titlesec}
\usepackage{xcolor}
\usepackage{fancyhdr}
\usepackage{array}
\usepackage{tabu}
\usepackage{hyperref}
\hypersetup{colorlinks=true, linkcolor=black, citecolor=black, urlcolor=gray}

\titleformat{\chapter}[hang] 
{\normalfont\LARGE\bfseries}{\thechapter}{1em}{}
\titlespacing*{\chapter}{0pt}{-20pt}{20pt}

\pagestyle{fancy}
\fancyhf{}
\fancyhead[LE,RO]{\thepage}
\fancyhead[RE]{\textit{\leftmark}}
\fancyhead[LO]{\textit{\rightmark}}
\renewcommand{\headrulewidth}{0.4pt}

\newenvironment{lorebox}{\begin{quote}\small\itshape}{\end{quote}}

\begin{document}

\frontmatter
\begin{titlepage}
\centering
\vspace*{2cm}
{\huge\bfseries Akilan\par}
\vspace{0.5cm}
{\Large The Southern Realms\par}
\vspace{2cm}
{\large A Fate's Edge Expansion\par}
\vspace{1.5cm}
\vfill
\end{titlepage}

\tableofcontents

\mainmatter

\chapter{Introduction: The Gulf of Oshiira}

South of the Kahfagian straits, beyond the Dolmis Narrows, the Amaranthine opens into a warmer, greener sea. This is the Gulf of Oshiira – a body of water that for centuries has been an Oshiiran lake. The great confederacy of the Sahel grasslands, the Oshiira Empire, controls its trade with an iron grip of logistics, treaties, and the legendary \textit{Widow’s Judgment}.

The **Shoreless Bay** marks Oshiira’s northern coastline – a broad, shallow sound where grain barges from the Sahel load for Ecktoria, Vhasia, and Silkstrand. Oshiiran wheat feeds half the Amaranthine; its dense canal cities are the breadbasket of the known world. Ecktoria and Oshiira fought four wars to a standstill, though the empire suffered worse. The memory of those wars lingers in every weigh‑house and every tariff.

Kahfagia was founded as a check on Ecktorian power – a Oshiiran colony granted independence to harry the empire’s western rival. Now it is a pirate‑republic, bound to Oshiira by a fragile non‑aggression pact and a cold war of spies, tariffs, and proxy raids.

To the east, the Ashaani magocracy nurses old wounds. Ashaan invaded Oshiira in the distant past and was broken. Under the Veiled Sisters, it has spent a century “recovering” – rebuilding its spirit‑fleets and dreaming of the Brass Gate. The desert that bears its name extends north into the Fharan Peninsula, linking the two realms by sand and salt.

To the south, the highlands of Taharka grow rich on monsoon trade; their horses are legends. Beyond them, the walking stone cities of Ngomebe move across the plateau, their Mason‑Kin Houses trading mineral wealth for Aeler engineering. In the deep jungle of Sekogo, druids whisper to spirits, and in the Crimson Valley, the Lethai‑ar guard their borders fiercely, their dark chestnut skin blending with the red earth.

This expansion provides cultural frameworks, travel decks, and a dozen new Patrons for the lands of southern Akilan. Use it to run campaigns of grain politics, slave‑liberation, stone‑city exploration, and the slow collapse of an empire that has grown too vast for its own ledgers.

\section*{Design Goals}
\begin{itemize}
\item \textbf{Empire as infrastructure:} Oshiira wins by counting bushels, not bodies.
\item \textbf{Magic as terrain:} Ashaan’s wizard‑lords warp the land; Sekogo’s druids \emph{are} the land.
\item \textbf{Stone as sovereignty:} Ngomebe does not build cities; it \emph{moves} them.
\item \strong{Cross‑cultural trade:} Akilan is a knot of hostile and hungry markets.
\item \textbf{Low bookkeeping, high tension:} Use clocks, Strings, SB menus as in other expansions.
\end{itemize}

\chapter{Peoples of Akilan}

\section{Oshiira – The Ledger Empire}

\textbf{Pillars:} Bureaucracy • Grain as power • The Widow’s judgment • Maritime tariffs

\textbf{Appearance:} Oshiira is a confederacy of hundreds of tribes and city‑states, united by the \textit{Grain Ledger}. The capital, Mbaro, is a city of canals, weigh‑houses, and the great Granary Tower. Oshiira is the breadbasket of the Amaranthine; its wheat fills the holds of Ecktorian, Vhasian, and Silkstrand ships. The Shoreless Bay is its northern border, a calm sound dotted with grain‑loading piers.

\textbf{History:} Oshiira fought four wars with Ecktoria, ending in a draw that left the empire bloodied and Ecktoria humiliated. The non‑aggression pact with Kahfagia holds only because a real war would sink both; instead, they wage a cold war of spies and tariffs.

\textbf{Voice:} An Oshiiran scribe sounds \textbf{haughty} – precise, condescending, and convinced that the ledger is the only truth. They speak of “bushel‑equivalents” and “tare weights” as if they were sacred texts.

\textbf{Etiquette Hooks:}
\begin{itemize}
\item \textbf{Present a stamped receipt} – gain Position +1 in petition scenes.
\item \textbf{Count aloud} – reduce DV by 1 for Broker tests.
\item \textbf{Declare your neutrality} – cancel one Oshiiran suspicion SB.
\item \textbf{Name the Widow first} – never speak of succession before her title.
\end{itemize}

\textbf{Strings:} Grain Ledger page, Tariff stamp, Barge priority token, Weigh‑house seal, Treaty seal (from the cold war pact).

\textbf{Factions:}
\begin{itemize}
\item \emph{The Vermilion Corps} – internal auditors, feared and hated.
\item \emph{The Canal Mothers} – matriarchs who control water‑rights.
\item \emph{The Salt Syndicate} – merchants who own the desert caravans.
\item \emph{The Narrow Captains} – naval officers who itch to provoke Kahfagia.
\end{itemize}

\section{Kahfagia – The Flag of Storms}

(Keep all original lore; only add the cold war dynamics.)

\textbf{Voice:} A Kahfagian speaks with an \textbf{educated Caribbean/Colonial} accent – sharp, musical, and laced with legalisms learned from Oshiiran clerks and Vilikari smugglers. They use double negatives for emphasis and pepper their speech with maritime slang.

\textbf{Cold War with Oshiira:} The non‑aggression pact is a sword of Damocles. The Narrow of Sorrows is the front line: patrol boats shadow each other, customs officials plant spies in weigh‑houses, and every captured pirate is a deniable asset. The \textit{Admiralty of Red Shoal} maintains a secret line to the Widow’s Vermilion Corps – a back channel for when proxy violence goes too far.

\textbf{Etiquette Hooks:}
\begin{itemize}
\item \textbf{Fly a flag} – never approach a port without a visible banner (Position +1).
\item \textbf{Name your ransom} – DV –1 for hostage negotiations.
\item \textbf{Share a cup of salt water} – cancels the first deception SB.
\item \textbf{Fly the flag of truce} (white with a black crab) – Position +1 when approaching Oshiiran ships; misuse earns a Black Mark.
\end{itemize}

\textbf{Strings:} Storm‑flag, Lantern code, Letter of marque, Cipher book (half‑key to Oshiiran dispatches).

\section{Ashaan – The Veiled Throne}

\textbf{Pillars:} Magocracy • Spirit slavery • The Veil • The Knife

\textbf{Voice:} Ashaani speak with \textbf{duplicitous} elegance. They offer compliments that are also threats, and their promises always have a hidden clause. A Veiled Sister’s whisper can sound like silk and cut like glass.

\textbf{History:} Ashaan invaded Oshiira in the ancient past and was broken. The empire has spent a century “recovering” under the Sisters – rebuilding its spirit‑fleets, binding new demons, and dreaming of the Brass Gate. The Red Desert is a wound left by that old war; it now links Ashaan to the Fharan Peninsula, connecting the two desert realms.

\textbf{Etiquette Hooks:}
\begin{itemize}
\item \textbf{Veil your face} – show only your eyes before the Sisters (Position +1).
\item \textbf{Speak in thirds} – offer two options and a hidden third (GM may grant a clue).
\item \textbf{Never name the dead} – breaking this earns a geas.
\end{itemize}

\textbf{Strings:} Spirit‑lamp (one‑use question), Brass seal (guest status), Veil‑clip (speak without being silenced).

\section{Ngomebe – The Stone Men}

\textbf{Pillars:} Moving cities • Mason‑Kin law • Ancestral stone • Anti‑magic wards

\textbf{Appearance:} The Ngomebe are dark‑skinned (the Aeler are also black‑skinned in the southern ranges; the two peoples share ancient ties). Their cities walk – colossal stone structures on rollers and sledges, dragged by teams of oxen, magic, and Aeler‑designed capstans. This is one of the best Aeler/Human partnership dynamics in the known world: Aeler engineers maintain the bearings; Ngomebe masons carve the keystones.

\textbf{Voice:} Ngomebe speak in low, measured tones – the voice of stone. They rarely hurry a sentence.

\textbf{Etiquette Hooks:}
\begin{itemize}
\item \textbf{Touch the city} – place your palm on the outer wall (Position +1).
\item \textbf{Speak of repair, not conquest} – DV –1 for arbitration.
\item \textbf{Never carve a name} – breaking this grants the GM 1 SB.
\end{itemize}

\textbf{Strings:} Keystone charter, Sledge‑wax, Stone‑song fragment.

\section{Taharka – The Monsoon Crown}

\textbf{Pillars:} Mountain cavalry • Rock‑chapels • Spice trade • Neutrality

\textbf{Wealth:} Taharka grows rich on monsoon trade – spices, coffee, and incense that travel east to the Hintara Ocean and west to the Gulf. The kingdom has never been conquered; its passes are a graveyard for Ashaani armies.

\textbf{Voice:} Taharkans speak formally, with elaborate greetings and a deep respect for lineage. They will praise your horse before they praise you.

\textbf{Etiquette Hooks:}
\begin{itemize}
\item \textbf{Remove your shoes} – Position +1 for petitions.
\item \textbf{Name your horse first} – DV –1 for negotiation.
\item \textbf{Never refuse a cup} – breaking this earns a 2‑segment Penalty clock.
\end{itemize}

\section{Sekogo – The Green Council}

\textbf{Pillars:} Druidic council • Honey‑orchards • Spirit‑kin • No slavery

\textbf{Location:} Deep jungle between Oshiira and the southern plateau. The Crimson Valley, home to the Lethai‑ar, lies at its eastern edge.

\textbf{Voice:} Sekogo speak in riddles and plant metaphors. They never say “no” directly; they say “the tree does not grow there.”

\textbf{Etiquette Hooks:}
\begin{itemize}
\item \textbf{Offer honey} – DV –1 for trade.
\item \textbf{Speak to the trees} – Position +1 for parley.
\item \textbf{Do not count} – using exact numbers offends; the GM may bank 1 SB.
\end{itemize}

\section{Lethai‑ar of the Crimson Valley}

\textbf{Appearance:} The Lethai‑ar in Akilan are chestnut‑brown (darker than the maple‑skinned northern elves). They guard the Crimson Valley – a jungle rift between Sekogo and Oshiira – with fierce neutrality. Their silk bridges span the red canyons; their web‑law keeps both empires from using the valley as a shortcut to war.

\textbf{Relationship with Neighbors:} They trade with Sekogo druids (honey for poison) and with Ngomebe (stone‑song for silk). Oshiira respects their borders because any invasion would bleed through the valley.

\textbf{Marks:} Their skin‑tattoos (Marks) are woven with red ochre and jungle sap; context keys involve the phases of the crimson moon.

\section{Aeler of the Southern Ranges}

\textbf{Appearance:} The Aeler of Akilan are black‑skinned, their hair often dusted with iron ore. They live in the mountain ranges south of Oshiira, adjacent to Ngomebe. Their partnership with the Stone Men is legendary: Aeler vent‑priors design the bearing systems for the walking cities; Ngomebe masons provide the keystone hearts. The two peoples intermarry in the border holds, and their guilds share a single underworld.

\textbf{Culture:} Same as northern Aeler (deep accounting, breath‑bonds, keystone rights), but with a southern accent: the “oven” is a furnace for smelting, not baking; the “grain surety” is paid in iron blooms.

\section{Amaria – The Iron Shore}

\textbf{Pillars:} Ecktorian colony • Iron mines • Legal hybrid • Isolated

\textbf{Location:} A small peninsula on the Gulf’s west coast, squeezed between Kahfagia and Oshiira. It remains neutral by paying tribute to both.

\textbf{Etiquette Hooks:} as previously written.

\chapter{Travel Decks for Akilan}

Use the standard Fate’s Edge travel procedure (draw one card from each suit: Spade=place, Heart=actor, Club=pressure, Diamond=leverage). Highest rank sets clock size (4/6/8/10). Each deck below replaces the regional decks from the core travel guide for Akilan.

\section{Oshiira (Sahel Grasslands)}

\begin{tabu} to \linewidth {|X|X|}
\hline
\multicolumn{2}{|c|}{\textbf{Spades – Places}} \\
\hline
2–5 & Mudflat weigh‑station; grain sacks piled under reed awnings. \\
6–10 & Granary tower with a cracked ledger door. \\
J–K & Canal junction; three barges, one pilot with a sealed writ. \\
A & The Widow’s weigh‑house – bronze scales, silent clerks. \\
\hline
\multicolumn{2}{|c|}{\textbf{Hearts – People}} \\
\hline
2–5 & Grain inspector with a brass stamp. \\
6–10 & Canal Mother bargaining over water rights. \\
J–K & Vermilion Corps auditor, ink on fingers. \\
A & A Narrow Captain nursing a grudge against Kahfagia. \\
\hline
\multicolumn{2}{|c|}{\textbf{Clubs – Complications}} \\
\hline
2–5 & Dust from the Ashaani desert coats the grain (Inspector halts the caravan). \\
6–10 & Locust swarm (Supply –1). \\
J–K & Cold war provocation – a Kahfagian “pirate” is sighted. \\
A & The Widow’s health fails; succession rumors spike. \\
\hline
\multicolumn{2}{|c|}{\textbf{Diamonds – Leverage}} \\
\hline
2–5 & Receipt for stolen grain – blackmail material. \\
6–10 & Barge priority token – skip one queue. \\
J–K & Treaty seal from the cold war pact – safe passage from both navies (once). \\
A & A page from the Grain Ledger – can be sold or ransomed. \\
\hline
\end{tabu}

\section{Ashaan (Red Desert & Fharan Link)}

\begin{tabu} to \linewidth {|X|X|}
\hline
\multicolumn{2}{|c|}{\textbf{Spades – Places}} \\
\hline
2–5 & Salt pan with a broken spirit‑cage. \\
6–10 & Black marble temple; steps stained red. \\
J–K & The Brass Gate – a fortified strait with spirit‑buoys. \\
A & A Veiled Sister’s sanctum, entrance hidden by a mirage. \\
\hline
\multicolumn{2}{|c|}{\textbf{Hearts – People}} \\
\hline
2–5 & Slave‑hunter (Black Hand) with a branding iron. \\
6–10 & Breath‑less agent (freed slave, now spy). \\
J–K & A Veiled Sister testing the party’s honesty. \\
A & An Ashaani magician who offers dangerous knowledge. \\
\hline
\multicolumn{2}{|c|}{\textbf{Clubs – Complications}} \\
\hline
2–5 & Spirit‑dust storm (Navigaton DV +1, mark Fatigue). \\
6–10 & Brass Gate closed for “inspection” – delay 2 segments. \\
J–K & A bound spirit breaks free; attacks the nearest living thing. \\
A & The Veil fails – spirit‑visions haunt all present. \\
\hline
\multicolumn{2}{|c|}{\textbf{Diamonds – Leverage}} \\
\hline
2–5 & A slave’s true name – can be used to free them (or bind them further). \\
6–10 & Spirit‑lamp – one question answered by a bound entity. \\
J–K & Map of the Fharan salt routes. \\
A & A sealed brass box containing a Veiled Sister’s secret. \\
\hline
\end{tabu}

\section{Ngomebe (Stone Lands)}

\begin{tabu} to \linewidth {|X|X|}
\hline
\multicolumn{2}{|c|}{\textbf{Spades – Places}} \\
\hline
2–5 & A city on the move – rollers groan, oxen low. \\
6–10 & Abandoned quarry; stone‑song echoes. \\
J–K & Mason‑Kin hall, walls carved with lineage. \\
A & The edge of a city’s path – the stone is warm. \\
\hline
\multicolumn{2}{|c|}{\textbf{Hearts – People}} \\
\hline
2–5 & Stone‑cutter whose city has left them behind. \\
6–10 & Mason‑Kin elder, grey with granite dust. \\
J–K & Root‑binder heretic, planting acorns in the rollers. \\
A & Aeler engineer: grease‑stained, carrying a level. \\
\hline
\multicolumn{2}{|c|}{\textbf{Clubs – Complications}} \\
\hline
2–5 & A keystone cracks – the city stops. \\
6–10 & Root‑binder sabotage – a wheel jams. \\
J–K & Mason‑Kin feud over the next migration route. \\
A & Earth tremor (all actions Desperate for a scene). \\
\hline
\multicolumn{2}{|c|}{\textbf{Diamonds – Leverage}} \\
\hline
2–5 & Keystone charter – right to dwell in a city for a season. \\
6–10 & Sledge‑wax – greases negotiation with the Mason‑Kin. \\
J–K & Stone‑song fragment – can calm a city’s spirit. \\
A & Aeler‑Ngomebe partnership seal – access to both guilds. \\
\hline
\end{tabu}

\section{Sekogo & Crimson Valley}

\begin{tabu} to \linewidth {|X|X|}
\hline
\multicolumn{2}{|c|}{\textbf{Spades – Places}} \\
\hline
2–5 & Honey‑orchard; the trees hum. \\
6–10 & Druid circle; stones covered in moss. \\
J–K & A Lethai‑ar silk bridge spanning a crimson gorge. \\
A & The Ungwe – a tree so old it speaks. \\
\hline
\multicolumn{2}{|c|}{\textbf{Hearts – People}} \\
\hline
2–5 & Sekogo honey‑gatherer with a fruit offering. \\
6–10 & Green Council druid carrying a thorn‑staff. \\
J–K & Waker radical plotting to awaken the jungle. \\
A & Lethai‑ar silk‑warden, masked and quiet. \\
\hline
\multicolumn{2}{|c|}{\textbf{Clubs – Complications}} \\
\hline
2–5 & A thorn‑spirit blocks the path – must offer honey or fight. \\
6–10 & The trees remember a broken promise. \\
J–K & Waker summons a fever‑mist (all actions Desperate). \\
A & Lethai‑ar close the silk bridge – no crossing until a debt is paid. \\
\hline
\multicolumn{2}{|c|}{\textbf{Diamonds – Leverage}} \\
\hline
2–5 & Honey‑charter – safe passage through one druid’s territory. \\
6–10 & Ungwe leaf – the tree’s blessing (one re‑roll on a Survival test). \\
J–K & A name spoken by the jungle – reveals a hidden path. \\
A & Lethai‑ar web‑token – parley under silk, once. \\
\hline
\end{tabu}

\section{Taharka (Highlands)}

\begin{tabu} to \linewidth {|X|X|}
\hline
\multicolumn{2}{|c|}{\textbf{Spades – Places}} \\
\hline
2–5 & Terraced farm; coffee shrubs on every step. \\
6–10 & Rock‑chapel; the walls drip condensation. \\
J–K & High pass; snow even in summer. \\
A & The King’s stable – horses with brass halter. \\
\hline
\multicolumn{2}{|c|}{\textbf{Hearts – People}} \\
\hline
2–5 & Horse trader with a copper bell. \\
6–10 & Ridge‑Watcher captain, bow over shoulder. \\
J–K & Spice House merchant, robes smelling of cardamom. \\
A & The Crown Council’s envoy – careful, neutral. \\
\hline
\multicolumn{2}{|c|}{\textbf{Clubs – Complications}} \\
\hline
2–5 & A landslide blocks the pass (delay 2 segments). \\
6–10 & Ashaani raiders disguised as merchants. \\
J–K & Horse sickness – all mounts lose 1 Speed. \\
A & The King is assassinated; pass closes to all foreign traffic. \\
\hline
\multicolumn{2}{|c|}{\textbf{Diamonds – Leverage}} \\
\hline
2–5 & Horse‑bond – a foal promised (future String). \\
6–10 & Pass‑stone – crossing a high pass at half toll. \\
J–K & Coffee bean tally – small bag, accepted as currency in any Taharkan market. \\
A & A secret treaty between Taharka and Oshiira – if exposed, it could end neutrality. \\
\hline
\end{tabu}

\section{Kahfagia (Cold War Theater)}

Use the core Kahfagia travel deck; add these cold‑war specific entries:

\textbf{Spades (Place) – The Narrow of Sorrows} (any rank): A reef‑laced channel where Oshiiran and Kahfagian warships watch each other across a cannon’s breadth.

\textbf{Hearts (Person) – The Narrow Captain} (any rank): Oshiiran naval officer who chafes at the pact; may offer a deal to provoke an incident.

\textbf{Clubs (Complication) – Flag of Truce Misused} (any rank): A Kahfagian privateer flies the white crab flag to ambush an Oshiiran grain barge. The party is caught in the middle.

\textbf{Diamonds (Leverage) – Cipher Book} (any rank): Half‑key to intercepted Oshiiran dispatches; one use to learn a secret.

\chapter{African-Inspired Patrons of Akilan}

These twelve Patrons are drawn from the cultures, landscapes, and conflicts of Akilan. Each follows the standard Fate’s Edge Patron format: Lore, Favored Deeds, Taboos, and Rites (Low/Standard/High). Obligation capacity = Spirit + Presence.

\section{Mbira – The First Weigh}

\textbf{Domain:} Grain, census, justice, the unalterable ledger.

\textbf{Lore:} Mbira is the spirit of the first weigh‑house, present when the first bushel of grain was measured against the first bronze weight. She does not forgive arithmetic errors. Her servants are auditors, judges, and tax collectors who believe that a correctly filed receipt is holier than a temple.

\textbf{Favored Deeds:} Correcting a false ledger, exposing a grain thief, paying a debt in full with interest.

\textbf{Taboos:} Burning a receipt, forgiving a debt without witness, using false weights.

\textbf{Low Rite – Stamp of Truth (4 XP):} Touch a document; for one scene, any forgery attempted on it becomes obvious to you. (Mark +1 Obligation)

\textbf{Standard Rite – Weigh the Heart (8 XP):} Target must confess a minor financial crime or suffer –1 die to all social rolls until they do. (Mark +2 Obligation)

\textbf{High Rite – The Ledger of Souls (14 XP):} For one hour, you can see the “value” of every person you meet – their debts, their credits, their unfulfilled contracts. (Mark +3 Obligation; backlash: you may become obsessed with balancing accounts)

\section{Tano – The Stone that Walks}

\textbf{Domain:} Migration, keystones, the pact between earth and engineering.

\textbf{Lore:} Tano is a spirit of the moving cities – the consciousness that allows a stone to roll without cracking. He speaks in the language of bearings and stresses. His followers are Aeler engineers and Ngomebe masons who believe that a well‑greased sledge is a prayer.

\textbf{Favored Deeds:} Repairing a city’s keystone, clearing a migration route, settling a Mason‑Kin dispute.

\textbf{Taboos:} Letting a city stop without cause, using a chisel on living rock, harming an ox that pulls a city.

\textbf{Low Rite – Song of the Keystone (4 XP):} Sing a low note; one stone door or bridge opens or closes. (Mark +1 Obligation)

\textbf{Standard Rite – Bearings of Peace (8 XP):} For one leg, your mount or wagon ignores the first difficult terrain penalty. (Mark +2 Obligation)

\textbf{High Rite – Tano’s Path (14 XP):} You may move a walking city up to one mile in a single night – but you must sacrifice an object of personal value to the rollers. (Mark +3 Obligation)

\section{Nkosi – The Lion of the Sahel}

\textbf{Domain:} Grasslands, courage, the pride of the hunt.

\textbf{Lore:} Nkosi is a great lion spirit who tests the strong and eats the weak. He respects only those who face him standing. His followers are warriors, scouts, and mercenaries who believe that retreat is permissible only if it sets a trap.

\textbf{Favored Deeds:} Defending the helpless, winning a duel against a stronger foe, sparing an enemy who surrenders.

\textbf{Taboos:} Ambushing from behind, harming a cub, breaking a sworn oath of hospitality.

\textbf{Low Rite – Mane of Fire (4 XP):} Your hair (or helmet) glows with golden light; +1 die to Intimidation for one scene. (Mark +1 Obligation)

\textbf{Standard Rite – Roar of the Pride (8 XP):} All allies within Near range gain +1 die to resist fear for one scene. (Mark +2 Obligation)

\textbf{High Rite – Nkosi’s Judgment (14 XP):} Challenge one opponent to a duel. If you win, they must perform one service for you (no suicide). If you lose, you owe them a season of fealty. (Mark +3 Obligation)

\section{Maji – The Mother of Wells}

\textbf{Domain:} Water, wells, the hollow spaces in stone.

\textbf{Lore:} Maji is the spirit of every spring, cistern, and rain‑catch. She is generous but jealous: if you waste her gift, she will dry your throat. Her followers are well‑diggers, canal mothers, and desert guides.

\textbf{Favored Deeds:} Digging a new well, sharing water with the thirsty, cleaning a cistern.

\textbf{Taboos:} Polluting a spring, hoarding water during a drought, lying near a well.

\textbf{Low Rite – Cup of Mercy (4 XP):} Pour water on dry ground; for one scene, you know the direction to the nearest clean water. (Mark +1 Obligation)

\textbf{Standard Rite – Seal the Cistern (8 XP):} Touch a water source; it becomes undrinkable to all but you and your allies for one day. (Mark +2 Obligation)

\textbf{High Rite – Maji’s Tears (14 XP):} Cause a short rain over a one‑mile area. Crops grow, but the rain also exposes any spirit within the area. (Mark +3 Obligation)

\section{Kalia – The Veiled Sister}

\textbf{Domain:} Secrets, spirit‑binding, the hidden knife.

\textbf{Lore:} Kalia is a rogue Veiled Sister who ascended to patronship by binding her own shadow. She teaches that every truth has a price and that the best secrets are the ones you never speak. Her followers are spies, assassins, and liberators of slaves.

\textbf{Favored Deeds:} Freeing a spirit from bondage, revealing a tyrant’s secret, wearing a veil in a foreign court.

\textbf{Taboos:} Breaking a vow of silence, binding a child’s spirit, speaking your true name to an enemy.

\textbf{Low Rite – Knife of Silence (4 XP):} For one scene, your footsteps make no sound. (Mark +1 Obligation)

\textbf{Standard Rite – Unbind (8 XP):} Release a bound spirit from an object; the spirit owes you one question. (Mark +2 Obligation)

\textbf{High Rite – Kalia’s Mask (14 XP):} For one day, you may appear as any person you have touched. The illusion breaks if you draw blood. (Mark +3 Obligation)

\section{Sefu – The Forge of Stone}

\textbf{Domain:} Iron, smithing, the fire under the plateau.

\textbf{Lore:} Sefu is a spirit of the great furnaces where Ngomebe and Aeler smelt their ore together. He is patient, hot, and unforgiving of impurities. His followers are smiths, miners, and metal‑traders.

\textbf{Favored Deeds:} Forging a weapon for a worthy cause, returning a stolen ingot, teaching a child to strike true.

\textbf{Taboos:} Selling flawed iron, using a forge for murder, allowing a fire to go unattended.

\textbf{Low Rite – Steady Hand (4 XP):} For one scene, your Craft rolls gain +1 die. (Mark +1 Obligation)

\textbf{Standard Rite – Sefu’s Seal (8 XP):} Touch a metal object; for one day, it cannot be melted or reshaped. (Mark +2 Obligation)

\textbf{High Rite – The Living Anvil (14 XP):} You may forge a single weapon that ignores mundane armor for one battle. At the end, the weapon crumbles. (Mark +3 Obligation)

\section{Zuri – The Honey‑Mother}

\textbf{Domain:} Bees, sweetness, the pact of pollination.

\textbf{Lore:} Zuri is a spirit of the Sekogo groves, who taught the druids how to speak to bees. She is kind, but she remembers every sting. Her followers are beekeepers, healers, and negotiators.

\textbf{Favored Deeds:} Planting a flowering tree, rescuing a swarm, sharing honey with the hungry.

\textbf{Taboos:} Burning a hive, using smoke to harvest, killing a bee without need.

\textbf{Low Rite – Sweet Tongue (4 XP):} For one scene, your Sway rolls gain +1 die. (Mark +1 Obligation)

\textbf{Standard Rite – Zuri’s Blessing (8 XP):} Touch a wound; it heals 1 Harm, but the patient craves honey for a day. (Mark +2 Obligation)

\textbf{High Rite – The Swarm’s Voice (14 XP):} For one hour, you may command an ordinary swarm of bees (treat as Cap 3 ally). (Mark +3 Obligation)

\section{Kito – The Rider of the Monsoon}

\textbf{Domain:} Horses, speed, the wind that brings rain.

\textbf{Lore:} Kito is a spirit of Taharkan horsemanship, born from the first mare that outran the monsoon. He is proud, fleet, and impatient with slow talk. His followers are messengers, scouts, and horse traders.

\textbf{Favored Deeds:} Racing fairly, cooling a lathered horse, sharing a pass with a stranger.

\textbf{Taboos:} Whipping a tired horse, selling a sick animal as healthy, blocking a mountain road.

\textbf{Low Rite – Kito’s Breath (4 XP):} Your mount gains +1 Speed for one leg. (Mark +1 Obligation)

\textbf{Standard Rite – The Rider’s Bond (8 XP):} For one scene, you and your mount share senses (you see what it sees). (Mark +2 Obligation)

\textbf{High Rite – Monsoon Hoof (14 XP):} You may ride through a sandstorm or duststorm without penalty for one scene. (Mark +3 Obligation)

\section{Ayo – The Drummer of the Gulf}

\textbf{Domain:} Tides, storms, the rhythm of wind and wave.

\textbf{Lore:} Ayo is a spirit of the Gulf of Oshiira, who drums the waves against the shore. She is wild, unpredictable, and generous to sailors who respect her beat. Her followers are pilots, fishermen, and storm‑flag captains.

\textbf{Favored Deeds:} Saving a drowning sailor, dancing during a storm, leaving a coin on the tide line.

\textbf{Taboos:} Cutting a ship’s drum, silencing a bell during a squall, fishing during a red tide.

\textbf{Low Rite – Drumbeat of Calm (4 XP):} For one scene, your ship ignores the first weather penalty. (Mark +1 Obligation)

\textbf{Standard Rite – Call the Tide (8 XP):} Raise or lower the tide in a small bay by one foot. (Mark +2 Obligation)

\textbf{High Rite – Ayo’s Dance (14 XP):} Summon a brief squall that scatters enemy ships (treat as Heist complication). (Mark +3 Obligation)

\section{Nuru – The Light of the Weigh‑house}

\textbf{Domain:} Lamps, ledgers, the clarity of numbers.

\textbf{Lore:} Nuru is a spirit of Oshiiran bureaucracy, who kindled the first lamp in the Granary Tower. He is precise, unforgiving, and utterly neutral. His followers are clerks, notaries, and the Vermilion Corps.

\textbf{Favored Deeds:} Reconciling a disputed account, illuminating a dark archive, catching a counterfeiter.

\textbf{Taboos:} Smudging a stamp, burning a receipt, lying under a lantern.

\textbf{Low Rite – Nuru’s Candle (4 XP):} Light a small flame; for one scene, you can see invisible ink. (Mark +1 Obligation)

\textbf{Standard Rite – The Auditing Gaze (8 XP):} For one scene, you can detect when someone is lying about numbers (weights, measures, counts). (Mark +2 Obligation)

\textbf{High Rite – The Ledger of Light (14 XP):} Touch a document; it becomes impossible to alter by any means for one year. (Mark +3 Obligation)

\section{Oya – The Wind of the Sahel}

\textbf{Domain:} Dust, change, the breath that clears the old.

\textbf{Lore:} Oya is a spirit of the harmattan – the dry wind that sweeps the grasslands, carrying red dust from Ashaan. She destroys the weak and cleanses the strong. Her followers are revolutionaries, deserters, and reformers.

\textbf{Favored Deeds:} Burning a corrupt contract, opening a sealed door, revealing a hidden chamber.

\textbf{Taboos:} Building a wall that blocks the wind, imprisoning a truth‑teller, hoarding grain during a drought.

\textbf{Low Rite – Oya’s Breath (4 XP):} For one scene, you may extinguish any non‑magical fire. (Mark +1 Obligation)

\textbf{Standard Rite – Dust Mask (8 XP):} For one scene, you are immune to the effects of sandstorms. (Mark +2 Obligation)

\textbf{High Rite – The Cleansing Wind (14 XP):} Summon a dust storm that lasts one hour. All actions within are Desperate. (Mark +3 Obligation)

\section{Ibeji – The Twin Stones}

\textbf{Domain:} Partnership, balance, the union of Aeler and Ngomebe.

\textbf{Lore:} Ibeji are twin spirits – one of stone, one of gear – who dance together to keep the walking cities alive. They teach that no craft is complete without another. Their followers are engineers and masons who work in pairs.

\textbf{Favored Deeds:} Building a bridge together, sharing credit for a repair, saving a partner’s life.

\textbf{Taboos:} Taking sole credit, letting a partner fail without help, breaking a twin‑keystone.

\textbf{Low Rite – Balanced Weight (4 XP):} For one scene, you and an ally gain +1 die to any roll made to assist each other. (Mark +1 Obligation)

\textbf{Standard Rite – Ibeji’s Seal (8 XP):} Touch two objects; for one day, they cannot be separated by less than magical force. (Mark +2 Obligation)

\textbf{High Rite – The Dance of Stone (14 XP):} You and an ally may move a stone city’s roller one mile without oxen – but you both mark 2 Fatigue. (Mark +3 Obligation)

\chapter{Campaign Fronts & Clocks}

Use these global clocks for an Akilan campaign.

\begin{tabu} to \linewidth {|X|c|c|c|X|}
\hline
\textbf{Faction} & \textbf{Influence} & \textbf{Stability} & \textbf{Crisis} & \textbf{Goal} \\
\hline
Oshiira Empire & 6 & 8 & 4 (Grain riots) & Maintain the Widow’s Peace \\
Ngomebe Mason‑Kin & 4 & 6 & 4 (City drift) & Keep the cities moving south \\
Ashaan Magocracy & 5 & 5 & 6 (Spirit revolt) & Capture the Brass Gate \\
Sekogo Green Council & 4 & 5 & 3 (Waker violence) & Protect the groves \\
Taharka Crown & 5 & 7 & 4 (Pass closure) & Maintain neutrality \\
Kahfagian Admiralty & 6 & 4 & 8 (Widow’s intervention) & Stay independent \\
Amarian Censorate & 4 & 4 & 5 (Mine strikes) & Extract silver \\
Lethai‑ar of the Crimson Valley & 3 & 6 & 2 (Silk bridges) & Guard the valley \\
\hline
\end{tabu}

\textbf{Front Clocks (campaign level):}
\begin{itemize}
\item \emph{The Widow’s Health} [8] – when filled, the empress dies; succession crisis triggers a region‑wide war.
\item \emph{The Spirit Debt} [6] – Ashaan’s bound spirits grow restless; when filled, a mass haunting destabilises the Veil.
\item \emph{The Stone Pact} [6] – Oshiira and Ngomebe negotiate a permanent border; when filled, either peace or war.
\item \emph{Cold War Pressure} [6] – ticks when the party raids an Oshiiran barge for Kahfagia, or smuggles grain past a Kahfagian blockade. At 2: sanctions (goods cost +1). At 4: a naval skirmish (Heist with Difficulty +1). At 6: the pact frays; open war becomes possible.
\end{itemize}

\chapter{Adventure Seeds}

Ten drop‑in scenarios for Akilan:

1. \textbf{The Widow’s Ledger:} A page from the Grain Ledger is stolen. Without it, a tribe cannot prove its harvest. Find the forger before the tribe starves.

2. \textbf{The City That Stopped:} A Ngomebe city has run aground. The Mason‑Kin blame a Root‑binder. The party must find the saboteur – or push the city themselves.

3. \textbf{The Spirit Cargo:} The party is hired to escort a sealed brass coffin from Ashaan to Oshiira. Inside is a bound spirit – and a rival faction wants to free it.

4. \textbf{The Honey Tithe:} A Sekogo grove demands tribute before letting a caravan pass. The tribute is a story. The party must tell a tale that satisfies the spirits.

5. \textbf{The Coffee Challenge:} A Taharkan prince offers the party a fortune to race his horse against an Oshiiran steed. The loser’s kingdom forfeits a pass.

6. \textbf{The Storm‑Flag Feud:} Two Kahfagian captains claim the same flag. The Admiralty demands the party referee a duel – to the first blood or the first confession.

7. \textbf{The Mine Strike:} Amarian miners barricade the shaft. The Censorate wants them removed. The party must negotiate – or break the barricade without causing a massacre.

8. \textbf{The Brass Gate:} An Ashaani magician has bribed the Oshiiran admiral of the straits. The party must expose the bribe before the Veiled Sisters sail through.

9. \textbf{The Walking Tree:} A Sekogo druid claims one of the Ngomebe cities crushed a sacred grove. The Mason‑Kin deny it. The party must survey the path – and find a witness that cannot lie.

10. \textbf{The Widow’s Heir:} The empress is dying. Her four children vie for the throne. The party is offered a barony to escort the heir of their choice to Mbaro. But each heir has a secret: one is a half‑breed, one is a spirit, one is a puppet, and one is already dead.

\chapter{GM Toolkit}

\section{Sahel Weather & Hazards (d12)}

\begin{enumerate}
\item Khamsin – dust obscures; ranged attacks –1 die.
\item Heat shimmer – mirages; Navigation DV +1.
\item Sudden flood – dry wadi becomes a river; lose 1 Supply.
\item Locust swarm – eat exposed grain; advance the Grain Ledger clock.
\item Desert night – cold enough to freeze water; mark 1 Fatigue if unsheltered.
\item Red dust – Ashaani sand carried eastward; social rolls –1 die; on a 1, a minor spirit manifests.
\item Sand geyser – pillar of sand; test Body + Athletics or be separated from the caravan.
\item Calm – no event; reroll.
\item Elephant stampede – test Stealth or take Harm 1 (trampling).
\item Grass fire – spread of flames; all actions Desperate until fire is passed.
\item Rain of frogs – omen; next scene begins with a 4‑clock “Sign from Maji”.
\item Widow’s Omen – distant thunder; the party dreams of ledgers and death.
\end{enumerate}

\section{Random Oshiiran Encounters (d12)}

\begin{enumerate}
\item Grain inspector demanding a bribe.
\item Vermilion Corps agent auditing a caravan.
\item Canal Mother seeking a new alliance.
\item Salt syndicate offering a loan (with interest).
\item Tribal chief asking for a witness to a feud.
\item Escaped Ashaani spirit, lost and confused.
\item Ngomebe stone‑cutter who has lost his city.
\item Kahfagian pirate selling “insurance”.
\item Amarian miner with a silver nugget and a story.
\item Taharkan horse‑trader with a mount that talks.
\item Sekogo druid asking for help against a Waker.
\item A Veiled Sister in disguise, testing the party’s honesty.
\end{enumerate}

\section{Integration with Existing Expansions}

\begin{itemize}
\item \textbf{Political Intrigue:} Oshiiran factions use Influence/Stability/Exposure. The Widow’s death is a Crisis trigger.
\item \textbf{Caravans:} Sahel legs require Water‑Wise and Grain‑Ledger tags. New Intent: Tribute.
\item \textbf{Amaranthine Sea:} The Gulf is a new theater; add Reef‑Singers and Monsoon Front.
\item \textbf{Witches of Fate's Edge:} Sekogo druids use Echo/Veil/Flow. The Veiled Sisters are a variant of the Cord (spirit‑binding instead of memory‑weaving).
\item \textbf{Horror Campaigns:} Ashaan is a Dread Clock generator (spirit slavery). Use the Reality Fracture track for the Veil.
\item \textbf{Shadows \& Steel:} Amaria and Kahfagia are criminal havens. Use Crew Heat and Informant Webs.
\item \textbf{Dragon’s Lair:} The walking cities of Ngomebe can be treated as lairs; the city itself is a boss.
\item \textbf{Kon’reh:} The Oshiiran school (Canray) is “Logistics Control”. Ashaani play “Veil \& Knife”. Sekogo refuse to play.
\end{itemize}

\chapter{Colophon}

\vspace{2\baselineskip}
\begin{lorebox}
Akilan is not a dark continent. It is a complex web of empires, spirits, and stone. The Widow’s Peace is a bargain, not a blessing. The cities move because staying still is death. And the spirits, the stones, and the trees all have opinions – if you know how to ask.

When the monsoon turns and the grain ships race the storms, you will see the ledger’s true weight. It is not written in ink. It is written in the bones of those who could not pay.

— \textit{Dusana of the Raven Road, at the weigh‑house of Mbaro, after a cup of coffee and a handful of salt.}
\end{lorebox}

\end{document}
